\section{Literaturrecherche 3D-Netze}
\label{sec:literaturrecherche 3D-Netze}

Niemeier (2008, S. 128) nennt zwei Wege, die Messungen mit dem Koordinatensystem gemeinsam durchführen. 
Die Reduktion auf ein Abbildungssystem (z. B. LV95, UTM) oder Formulierung der Beobachtungsgleichungen direkt in einem räumlichen System. 
Für hochgenaue lokale Ingenieurnetze ist ein 3D-System von Bedeutung, wenn die Höhe streng mitverarbeitet werden soll. 
Dafür die Kombination von GNSS mit terrestrischen Messungen ist das globale kartesische System gebräuchlich.

\subsection{Datenaufbereitung}
\label{sec:datenaufbereitung}

\subsection{Datenauswertung}
\label{sec:datenauswertung}

\subsection{Resultate}
\label{sec:resultate}

\subsection{Diskussion der Resultate}
\label{sec:diskussion der resultate}